
\chapter{Introduction}\label{ch:introduction}

\section*{Motivation}

Modern reinforcement learning systems increasingly operate in settings where
experience data is distributed across multiple agents, devices, or geographic
locations. Consider a fleet of autonomous vehicles that must collectively learn
safe driving policies, or a network of industrial robots sharing a manufacturing
task while operating in different physical environments. In such scenarios,
centralising all data for training is not only logistically expensive but may
be legally or ethically infeasible due to data privacy requirements.

\emph{Federated Learning}~\cite{mcmahan2017fedavg} addresses this problem by
allowing participants to collaboratively train a shared model while keeping their
local data private. Rather than sending raw data to a central server, each
participant computes a local model update and transmits only that update. The
server aggregates the updates and broadcasts an improved global model. This
cycle continues until convergence.

\emph{Federated Reinforcement Learning} (FedRL) applies this paradigm to the
reinforcement learning setting, where each participant is an agent interacting
with its own environment. Agents compute local policy gradient estimates from
their collected trajectories and send these gradients to a coordination server.
The server aggregates the gradients and updates the shared policy. This approach
combines the sample-efficiency benefits of sharing experience across agents with
the privacy and scalability benefits of federated communication.

\section*{The Reproducibility Problem}

Despite growing theoretical interest in FedRL, the field suffers from a severe
lack of standardised tooling. Published algorithms are typically evaluated in
custom experimental setups against custom baselines. Reproducing results from a
FedRL paper often requires reimplementing the entire framework from scratch --- a
process that is time-consuming, error-prone, and frequently leads to unintentional
discrepancies. The absence of shared infrastructure makes fair comparison across
methods nearly impossible and slows the pace of scientific progress.

This problem is well-recognised in the supervised federated learning community,
which has responded with frameworks such as Flower~\cite{beutel2020flower},
PySyft, and FedML. These tools provide standardised communication protocols,
client-server abstractions, and experiment management. However, their support for
reinforcement learning workloads --- where the fundamental data unit is a
variable-length trajectory rather than a fixed-size minibatch --- remains limited
and largely unevaluated.

The algorithm this work is most directly concerned with is
FedPG-BR~\cite{flint2023fedpgbr}, a Byzantine-robust federated policy gradient
method. FedPG-BR extends the basic federated policy gradient approach with a
probabilistic filter that identifies and excludes Byzantine agents --- workers
who send malicious gradient updates to corrupt the shared policy. While the
theoretical contributions of FedPG-BR are significant, no open reference
implementation existed that would allow other researchers to reproduce its
results, compare against it fairly, or build upon it as a baseline.

\section*{Our Contribution}

This thesis presents the \textbf{Flower FRL Benchmark}, an open-source testbench
for federated policy gradient research. Built on top of the Flower federated
learning framework~\cite{beutel2020flower}, the benchmark offers:

\begin{itemize}
    \item A complete implementation of FedPG-BR~\cite{flint2023fedpgbr},
          including its SCSG variance reduction and probabilistic Byzantine filter.
    \item Two additional aggregation baselines: SVRPG~\cite{papini2018svrpg}
          (variance reduction without Byzantine filtering) and GOMDP (plain
          gradient averaging).
    \item A \emph{strategy plugin system} allowing researchers to register a
          custom aggregation algorithm with a single Python decorator and have it
          immediately available for comparison, without modifying any framework code.
    \item Simulation of seven Byzantine attack types from the FedPG-BR paper and
          related literature.
    \item Out-of-the-box support for discrete control environments (CartPole-v1,
          Acrobot-v1, MountainCar-v0) and continuous control environments
          (LunarLander, HalfCheetah via MuJoCo~\cite{todorov2012mujoco}).
    \item A real-time web dashboard for experiment configuration, live reward
          monitoring, and per-agent diagnostics.
    \item Docker-based deployment for reproducible execution across platforms.
\end{itemize}

The overarching design goal is that a researcher should be able to implement and
test a new FedRL aggregation strategy in a single Python file, with no knowledge
of the underlying Flower communication protocol. The benchmark thus serves both
as a reproducible reference implementation of FedPG-BR and as an extensible
platform for future FedRL research.

\section*{Thesis Outline}

This thesis is structured as follows:

\begin{itemize}
    \item \textbf{\Cref{ch:background} (Background and Related Work):} We review
          the foundational concepts in reinforcement learning, policy gradient
          methods, federated learning, and Byzantine robustness, and describe the
          FedPG-BR and SVRPG algorithms in detail.
    \item \textbf{\Cref{ch:system} (System Design and Implementation):} We
          describe the architecture of the Flower FRL Benchmark, including its
          Flower integration, strategy plugin system, Byzantine attack simulation,
          and web dashboard.
    \item \textbf{\Cref{ch:conclusion} (Conclusion):} We summarise our
          contributions, discuss the limitations of the current implementation,
          and outline directions for future work.
\end{itemize}
